\documentclass[12pt]{article}

\title{A compact proof that choice from well-ordered partitions fails in NF, following Elliot Glazer}

\author{Lavinia Randall Holmes}

\usepackage{amssymb}

\begin{document}

\maketitle

This is a working note and could easily have mistakes in it.  The results in it are derived from work originally done by Elliot Glazer and ChatGPT Sol, his AI sidekick.  I have used my own expertise to give a faster way to prove that it is false in NF:  the combinatorics are lifted entirely from Glazer, and the final argument in terms of functions not increasing faster than T in the ordinals is mine (with a debt ultimately to Specker and C. Ward Henson).

The axiom of well-ordered choice (AC$_{\tt wo}$) is the assertion that every well-ordered partition has a choice set.
Given an arbitrary ordinal indexed sequence of sets $X_\alpha$, we can construct a sequence of disjoint
sets $Y_{T^3\alpha}=\iota^3``X_\alpha \times \{\alpha\}$ (if $s$ is a well-ordering, we define $s_{\alpha}$ as the item in the well-ordering with segment before it of order type $\alpha$: $s_\alpha$ is of type one lower than that of the segment, which is of type one lower than the ordinal index $\alpha$).
A choice set from the well-ordered partition obtained by dropping empty sets from the range of this sequence readily gives us a well-orderable set which meets each nonempty set in the original sequence.

For any cardinal $\kappa$, we define $\exp(\kappa)$ as the image under $T^{-1}$ of the cardinality of the power set of any  $X\in \kappa$.  We define $\aleph(\kappa)$ as the image under $T^{-2}$ of the
cardinality of the set of order types of well-orderings of subsets of any $X \in \kappa$.  These are precise analogues of cardinal operations in ordinary set theory and have familiar properties mod stratification considerations.  We do point out that we are using a type level ordered pair:  otherwise the exponent in the definition of $\aleph$ would be different.

We compute a bound on $\aleph(\exp(\kappa))$ for any cardinal $\kappa$.   Suppose that $X \in \kappa$.
Let $Y \subseteq {\cal P}(X)$ have a well-ordering.   Then $Y \times Y$ has a well-ordering.  Let $C_\alpha = (A_\alpha,B_\alpha)$ be an ordinal indexing of $Y\times Y$.  We consider the ordinal indexed sequence $D_\alpha=A_\alpha-B_\alpha$.  By AC$_{\tt wo}$ there is a well-orderable set $E$ which meets each nonempty $D_\alpha$.  The map $(A \in Y \mapsto A \cap E)$ is an injection, because if $A$ and $B$ are distinct elements of $Y$, there is an element of $A \Delta B$ in $E$ by construction.  We see that any well-ordering of a set of subsets of $X$ is isomorphic to a well-ordering of subsets of a well-ordered subset $E$ of $X$:  the cardinality of $E$ is bounded by $\aleph(\kappa)$.
So $\aleph(\exp(\kappa)) \leq \aleph(\exp(\aleph(\kappa)))$.

This implies further that $\aleph(\exp(\kappa)) < \aleph(\aleph(\exp(\aleph(\kappa))))$.  

For sufficiently large concrete $n$ (large enough that all computations on both sides succeed) we have 

$\aleph(T^{n-1}(|V|))=\aleph(\exp(T^n(|V|))) < \aleph(\aleph(\exp(\aleph(T^n(|V|)))))$. 

And this is impossible, for reasons which can be stated in very general terms.

Define a special function on cardinals of well-ordered sets as a function $F$ which commutes with $T$, is nondecreasing, and satisfies $F(\alpha)>\alpha$.

For no special function can we have $F(T(\alpha)) > \alpha$, if $T(\alpha)<\alpha$ and $\alpha$ is not dominated by a cantorian $\beta=T(\beta)$.

We prove this.  Suppose otherwise for a particular special $F$ and cardinal $\alpha$ of a well-ordered set.

Let $\gamma$ be smallest such that for some $n$, $F^n(\gamma) \geq \alpha$.

Note that $T(\gamma)$ is smallest such that for some $n$, $F^n(T(\gamma)) \geq T(\alpha)$, and by hypothesis this is the same cardinal:
clearly $F^{n+1}(T(\gamma)) \geq \alpha$ so $T(\gamma)\geq \gamma$, and we have $m$ such that $F^m(\gamma) \geq \alpha >T(\alpha)$, so $T(\gamma) \leq \gamma$.

Now let $N$ be smallest such that $F^N(\gamma)\geq \alpha$.  $N >0$ since $F^0(\gamma) = \gamma$, being cantorian, is certainly less than $\alpha$.

Certainly $T(N)$ is smallest such that $F^{T(N)}(\gamma) \geq T(\alpha)$.

We have $F^{N-1}(\gamma)<\alpha$ and $F^N(\alpha) \geq \alpha$.

We have $F^{T(N)-1}(\gamma)<T(\alpha)<\alpha$, $F^{T(N)}(\gamma) \leq T(\alpha)$, and by hypothesis $F^{T(N)+1}(\gamma) \leq \alpha$.

But this means that $N$ must be either $T(N)$ or $T(N)+1$.

We cannot have $N=T(N)$ because this would make $F^{TN}(\gamma) = F^N(\gamma)$ cantorian and greater than $\alpha$.

We cannot have $N=T(N)+1$ because $N$ and $T(N)$ must have the same parity (odd or even).

$F(\alpha) = \aleph(\aleph(\exp(\alpha)))$ is a special function on cardinals of well-ordered sets.

$\alpha = \aleph(T^{n-1}(|V|))$ satisfies $T(\alpha)<\alpha$ and is not dominated by any cantorian cardinal of well ordered sets.

So we have arrived at a contradiction.  So our original assumption of choice from well-ordered partitions is inconsistent with NF.

Note that this is a disproof of the Axiom of Choice itself, and might reveal more about what is really going on than the original proof.



\end{document}